\subsection{\textit{Sprint} 3}

Para esta \textit{sprint}, o time havia previsto finalizar a tela \textit{home} no \textit{front-end}, entregar
a tela de empresas, finalizar no \textit{back-end} todas as rotas básicas, que incluíam
\ac{cnae}s, municípios, estados e estabelecimentos, permitindo apresentar o \textit{front-end}
sem dados simulados. Além disso, o banco de dados precisava ser finalmente
definido com os \textit{scripts}, uma versão básica de pesquisa sobre presença digital
deveria ser adicionada e o time de \ac{ci-cd} precisava fazer o \textit{deploy} para a \ac{aws}
\cite{aws-docs}. Da
minha parte, estava previsto trabalhar em boa parte dos \textit{endpoints} do \textit{back-end} e
dar uma ajuda no \textit{front-end}, que, segundo comunicação com os colegas, estava
precisando de maior força-tarefa.

Da parte do \textit{front-end}, a equipe conseguiu entregar o prometido, inclusive já
começando a página de mapas, da qual participei no desenvolvimento. O \textit{deploy} de
\ac{ci-cd} se estendeu por muito tempo nesta \textit{sprint} e demandou um número grande de
recursos humanos, mas foi possível entregar a tempo o prometido. O banco de
dados passou a contar com os \textit{scripts} corretos, possibilitando rodar e acessar os
dados por meio do \textit{back-end}. Além disso, a equipe ficou confiante de que, dali em
diante, faríamos mínimas ou até nenhuma mudança nessa parte, que também se
estendeu um pouco mais do que o esperado.

O \textit{back-end}, como previsto, foi a parte em que mais participei. Atuei ativamente na
criação dos \textit{endpoints} de estados e \ac{cnae}s, inclusive aprendendo e criando testes
unitários para eles. Além disso, mediante observação de que alguns AGES I pareciam
perdidos e com poucas horas feitas, ajudei-os a se fazerem mais presentes na
equipe, marcando horários para \textit{pair programming} de \textit{endpoints} do \textit{back-end}. Nessas
ocasiões, transmiti conhecimento e respondi dúvidas relacionadas à organização de
pastas e boas práticas.

Esta foi a primeira \textit{sprint} deste projeto em que senti confiança para dizer que
desenvolvi por muito tempo sem encontrar grandes problemas no caminho. Acredito
que o mais próximo de um problema foi o fato de o banco de dados ainda estar meio
complicado de rodar, dificultando a realização de testes manuais. No entanto, esse
problema foi contornável a partir do momento em que começamos a desenvolver
testes automatizados, capazes de gerar dados simulados com o intuito de simular uma
entrada vinda do banco de dados.

Entrando na \textit{sprint} final, acredito que tínhamos um bom trabalho a ser entregue.
Naquele momento, pensei em me juntar com a equipe que estava desenvolvendo a
parte da presença digital e ajudar, visto que o cliente alegava que essa \textit{feature} era
a própria razão de existência do projeto e, portanto, parecia consideravelmente
atrasada em comparação às expectativas. Por fim, com o restante do tempo,
planejei continuar ajudando os AGES I com o desenvolvimento de \textit{endpoints} e,
dependendo do andamento das primeiras tarefas, implementar testes unitários no
projeto.
